\documentclass[11pt,a4paper]{article}

\usepackage[utf8]{inputenc}
\usepackage[T1]{fontenc}
\usepackage[spanish,es-nodecimaldot]{babel}
\usepackage[margin=2cm]{geometry}
\usepackage{lmodern}
\usepackage{microtype}
\usepackage{enumitem}
\usepackage[hidelinks]{hyperref}
\usepackage{titlesec}

\titlespacing*{\section}{0pt}{8pt}{3pt}
\setlist{nosep, leftmargin=*}
\setlength{\parskip}{3pt}
\setlength{\parindent}{0pt}

\title{\vspace{-1.2cm}\textbf{Detección de exoplanetas en curvas de luz de Kepler:\\
features físicas versus redes convolucionales 1D}\\[2pt]
\large INF396 -- Propuesta de proyecto}
\author{Noelia Morales \quad Matías Lara\\
\small Ingeniería Civil Informática, Universidad Técnica Federico Santa María}
\date{\vspace{-0.6cm}}

\begin{document}
\maketitle

\section{Descripción del problema}

Este trabajo trata sobre la \textbf{detección automática de estrellas con
exoplanetas} a partir de sus \textbf{curvas de luz}, es decir, series de tiempo
del flujo luminoso medido por el telescopio espacial Kepler. Cuando un planeta
pasa frente a su estrella (\emph{tránsito}), bloquea una pequeña fracción de su
luz y produce una caída breve y periódica del flujo. Detectar esa señal es
difícil: la caída suele ser muy pequeña comparada con el ruido instrumental y
con la variabilidad propia de la estrella, y las estrellas con planetas
confirmados son extremadamente escasas.

En la literatura, las redes convolucionales han mostrado ser efectivas para
clasificar señales de tránsito en datos de Kepler~\cite{shallue2018,ansdell2018},
y los clasificadores de \emph{gradient boosting} sobre features extraídas de la
curva alcanzan resultados comparables con un costo computacional mucho
menor~\cite{malik2022}. Sobre el mismo dataset que usaremos, Priyadarshini y
Puri~\cite{priyadarshini2021} reportan un 99{,}62\% de \emph{accuracy} con un
ensamble de CNN. Sin embargo, como solo el 0{,}9\% del conjunto de prueba es
positivo, un clasificador que prediga ``sin planeta'' para todas las estrellas
ya obtiene un 99{,}12\%. Esto muestra que la \emph{accuracy} no sirve para
evaluar el problema y motiva un protocolo de evaluación más riguroso.

\section{Objetivo del trabajo}

\textbf{Nuestro objetivo es comparar un clasificador basado en features con
significado físico contra una red neuronal convolucional 1D para detectar
exoplanetas en curvas de luz de Kepler, con un protocolo de evaluación adecuado
al fuerte desbalance de clases, y analizar mediante técnicas de
interpretabilidad qué patrones aprende la red.} En concreto:
\begin{enumerate}
  \item \textbf{Baseline con features físicas:} además de estadísticos
  descriptivos, extraer features del espectro de Fourier (FFT) y del
  periodograma \emph{Box Least Squares} (BLS)~\cite{kovacs2002}: período,
  profundidad, duración y potencia del tránsito más probable. Entrenar sobre
  ellas Random Forest y XGBoost.
  \item \textbf{Modelo secuencial:} entrenar una CNN 1D directamente sobre la
  curva preprocesada (\emph{detrending} con filtro de mediana y normalización),
  manejando el desbalance con ponderación de clases y \emph{data augmentation}
  propia de series de tiempo (desplazamientos circulares, inversión temporal,
  ruido gaussiano).
  \item \textbf{Estudio de ablación:} medir el aporte de cada decisión de
  diseño (preprocesamiento, ponderación de clases, \emph{augmentation}, tipo de
  features) activándolas una a la vez.
  \item \textbf{Interpretabilidad:} aplicar mapas de saliencia y Grad-CAM
  1D~\cite{selvaraju2017} para verificar si la CNN basa sus predicciones en los
  tránsitos o en artefactos, y analizar los falsos positivos y negativos.
\end{enumerate}
Como métricas usaremos \emph{recall}, precisión, F1 y el área bajo la curva
precisión--\emph{recall} (PR-AUC). Dado que el conjunto de prueba original tiene
solo 5 positivos, las estimaremos con validación cruzada estratificada repetida
e intervalos de confianza por \emph{bootstrap}, además de reportarlas sobre el
conjunto de prueba oficial.

\section{Punto de partida}

El proyecto no continúa un trabajo anterior. Existen implementaciones públicas
relacionadas: el código de AstroNet~\cite{shallue2018}, que trabaja con curvas
ya plegadas en fase del catálogo DR24 y no con este dataset, y numerosos
notebooks en Kaggle que entrenan CNN o SVM sobre este mismo conjunto. Estos
últimos suelen reportar \emph{accuracy} sobre un conjunto de prueba con 5
positivos y aplican técnicas como SMOTE directamente sobre la serie cruda.
Nuestro aporte está en:
(i)~usar features físicas (BLS/FFT) en lugar de estadísticos genéricos;
(ii)~un protocolo de evaluación robusto al desbalance;
(iii)~un estudio de ablación sistemático; y
(iv)~el análisis de interpretabilidad de la CNN.

\section{Dataset}

Usaremos \emph{Exoplanet Hunting in Deep Space -- Kepler labelled time series
data}\footnote{\url{https://www.kaggle.com/datasets/keplersmachines/kepler-labelled-time-series-data}}
(licencia CC0). Contiene curvas de luz ya limpiadas por NASA, con 3.197
mediciones de flujo por estrella y una etiqueta binaria (planeta confirmado o
no). Más del 99\% de los datos proviene de la Campaña~3; para aumentar los
positivos, se añadieron estrellas con planetas confirmados de otras campañas.
\begin{itemize}[label=\textbullet]
  \item \texttt{exoTrain.csv}: 5.087 estrellas, de las cuales 37 tienen
  exoplanetas (0{,}73\%).
  \item \texttt{exoTest.csv}: 570 estrellas, de las cuales 5 tienen
  exoplanetas (0{,}88\%).
\end{itemize}
El dataset no incluye marcas de tiempo. Para calcular el BLS supondremos una
cadencia uniforme de aproximadamente 30 minutos (la cadencia larga de Kepler) y
evaluaremos la sensibilidad de los resultados a ese supuesto.

\section{Artículos relacionados}

\begingroup
\renewcommand{\section}[2]{}%
\small
\bibliographystyle{plain}
\bibliography{refs}
\endgroup

\end{document}
